% sec_01_surface_law
\section{Setup and the master surface law}\label{sec:surface-law}

AEGIS reduces dosimetry to a single scalar field on the body surface: the
absorbed power density $\Sab(\rr)$, in watts per square metre, deposited at a
skin point $\rr$. Every fidelity level in this report computes $\Sab$, and
every refinement (Fresnel transmission, polarisation, curvature, diffraction,
inter-body coupling, coherent beamforming) is a correction to one factor of a
single master law. This section states that law, defines each symbol, and
fixes the sign and time conventions that the diffraction and coherent sections
depend on.

\subsection{Geometric setup}\label{sec:setup-geometry}

Consider a monochromatic plane wave of angular frequency $\omega = 2\pi f$
propagating in the unit direction $\khat$, with time-averaged Poynting vector
$\mathbf{S}_{\mathrm{inc}} = \Sinc\,\khat$. The scalar $\Sinc$ is the incident
power spectral density (the magnitude of the time-averaged Poynting flux,
\si{\watt\per\square\metre}). The body surface $\Sigma$ is described by an
outward unit normal field $\nhat(\rr)$ and is assumed locally flat on the scale
of the wavelength, so that a Fresnel half-space model applies pointwise. This
local-flatness assumption is the one relaxed by the curvature and diffraction
refinements of \cref{sec:grid} and the later geometry sections.

\begin{definition}[Incidence cosine]\label{def:mu}
At a surface point $\rr$ with outward normal $\nhat(\rr)$, the incidence cosine
is
\begin{equation}\label{eq:mu-def}
  \mu(\rr) \;\equiv\; \nhat(\rr)\cdot(-\khat) \;=\; \cos\theta_i(\rr),
\end{equation}
where $\theta_i$ is the local angle of incidence. A front-facing point has
$\mu > 0$, a point in geometric shadow has $\mu \le 0$.
\end{definition}

\subsection{Time-averaged power balance}\label{sec:power-balance}

The inward power flux through a surface element $\diff A$ at $\rr$ is
$\diff P_{\mathrm{in}} = \Sinc\,\mu\,\diff A$: the beam carries $\Sinc$ per unit
area normal to $\khat$, and a tilted element of its own area $\diff A$
intercepts only the projected fraction $\mu = \cos\theta_i$. The specularly
reflected wave leaves at the same angle with power density
$\abs{r(\theta_i)}^2\,\Sinc$, where $r(\theta_i)$ is the Fresnel amplitude
reflection coefficient for the relevant polarisation, so the outward reflected
flux is $\diff P_{\mathrm{ref}} = \abs{r}^2\,\Sinc\,\mu\,\diff A$.

At millimetre-wave frequencies the skin depth is $\delta \sim \SIrange{0.3}{0.5}{\milli\metre}$,
far below any body curvature radius, so all transmitted power is absorbed in
the surface layer. Energy conservation at the interface gives
$\diff P_{\mathrm{abs}} = \diff P_{\mathrm{in}} - \diff P_{\mathrm{ref}}
= \Sinc\,\mu\,(1-\abs{r}^2)\,\diff A$.

\begin{definition}[Power transmission coefficient]\label{def:T}
For a given polarisation the power-absorption (transmission) coefficient is
$T(\theta_i) \equiv 1 - \abs{r(\theta_i)}^2$. Its normal-incidence value is
written $T_0 \equiv T(0)$.
\end{definition}

\subsection{The master surface law}\label{sec:master-law}

Dividing the absorbed flux by $\diff A$ and gating out back-facing points
yields the law that organises the entire report.

\begin{proposition}[Master surface law]\label{prop:master-law}
The absorbed power density at a surface point $\rr$ is
\begin{equation}\label{eq:master-law}
  \boxed{\;\Sab(\rr) \;=\; \Sinc\;T_0\;\ReLU\!\bigl[\nhat(\rr)\cdot(-\khat)\bigr]\;
  O(\rr,\khat)\;}
\end{equation}
where $\ReLU(\mu)=\max(0,\mu)$ and $O(\rr,\khat)\in\{0,1\}$ in the
geometric-optics limit is the visibility factor, smoothed by the diffraction gate
in later sections. Each factor is independent and is refined separately in later
sections.
\end{proposition}

The five factors carry the whole physics:
\begin{itemize}[leftmargin=1.4em,itemsep=1pt]
  \item $\Sinc$: incident power spectral density at $\rr$
    (\si{\watt\per\square\metre}), supplied by the ray tracer or environment
    model.
  \item $T_0$: normal-incidence power transmission into tissue
    (\cref{def:T}), the material factor. \Cref{sec:tissue} derives it from the
    tissue refractive index and \cref{sec:fresnel} generalises it to the
    angle- and polarisation-dependent $\Teff(\rr)$.
  \item $\ReLU[\mu]$: the cosine projection combined with a back-face cull. It
    is zero in geometric shadow and equal to $\cos\theta_i$ on the lit side.
    \Cref{sec:grid} and the diffraction sections soften this hard gate into a
    physically derived smooth transition near the shadow boundary.
  \item $O(\rr,\khat)\in\{0,1\}$ in the geometric-optics limit: the visibility or occlusion factor, unity for a
    directly illuminated point and reduced by self-shadowing, inter-body
    occlusion, and ambient occlusion. It is developed in the geometry and
    inter-body sections.
\end{itemize}

For unpolarised or circularly polarised illumination the material factor is the
scalar $T_0$ shown in \cref{eq:master-law}. The exact, polarisation-resolved
form replaces $T_0$ with the effective transmission $\Teff(\rr)$ of
\cref{sec:pol}, giving
\begin{equation}\label{eq:master-law-exact}
  \Sab(\rr) \;=\; \Sinc\;\Teff(\rr)\;\ReLU\!\bigl[\mu(\rr)\bigr]\;O(\rr,\khat),
\end{equation}
which is exact wherever the surface is locally flat. The coherent levels of
\cref{sec:grid} replace the incoherent power product entirely with a
quadratic form in the complex precoder, recovering \cref{eq:master-law-exact}
in the incoherent limit.

\subsection{Sign and time conventions}\label{sec:conventions}

The following convention is load-bearing. The Fresnel branch cuts, the
penetration depth, and the Fock attenuation in the diffraction sections all
depend on it, so we state it once and use it everywhere.

\begin{remark}[Refractive-index and time convention]\label{rem:convention}
Fields vary in time as $e^{+\mathrm{i}\omega t}$. With this convention a lossy
medium has complex relative permittivity
$\epstilde = \varepsilon_r - \mathrm{i}\,\sigma/(\omega\varepsilon_0)$ and
complex refractive index
\begin{equation}\label{eq:n-convention}
  \ntilde \;=\; \sqrt{\epstilde} \;=\; n - \mathrm{i}\kappa,
  \qquad n>0,\;\kappa>0,
\end{equation}
where the branch is fixed by $\RE(\ntilde)>0$ and $\IM(\ntilde)<0$ so that a
wave decays as it penetrates the tissue. Throughout, $\mu = \nhat\cdot(-\khat)$
denotes the incidence cosine of \cref{def:mu}. Adopting the opposite
($e^{-\mathrm{i}\omega t}$) convention flips the sign of $\kappa$ and of every
imaginary part below.
\end{remark}
